\section{Введение}

\begin{frame}
    \frametitle{Содержание}
    \begin{enumerate}
        \item Введение
        \item Глава 1. Методика построения и аппроксимации многообразия решений
        \item Глава 2. Плоская круговая модель и конечные формулы
        \item Глава 3. Линейные \(R\)-преобразования
        \item Глава 4. Применение в баллистическом проектировании
        \item Заключение
    \end{enumerate}
\end{frame}

\begin{frame}
    \frametitle{Актуальность}
    \begin{columns}[T,onlytextwidth]
        \begin{column}{0.46\textwidth}
            \begin{itemize}
                \item Электрореактивные двигательные установки расширяют возможности малых межпланетных аппаратов.
                \item Высокий удельный импульс приводит к длительным управляемым перелётам и усложняет баллистическое проектирование.
                \item На ранних этапах анализа требуется быстро строить начальные приближения для последующего решения краевых задач.
            \end{itemize}
        \end{column}
        \begin{column}{0.5\textwidth}
            \centering
            \includegraphics[width=\linewidth]{chap4_mars_cleanver_rebuilt_Oxy}

            \vspace{0.15cm}
            \footnotesize Начальное приближение и уточнённая траектория для межпланетного перелёта
        \end{column}
    \end{columns}
\end{frame}

\begin{frame}
    \frametitle{Проблема проектирования класса перелётов}
    \begin{columns}[T,onlytextwidth]
        \begin{column}{0.54\textwidth}
            \centering
            \includegraphics[width=0.9\linewidth]{Startdate}
        \end{column}
        \begin{column}{0.42\textwidth}
            \small
            \textbf{Численное решение одной КЗ}
            \begin{itemize}
                \item фиксированы дата старта и целевая орбита;
                \item начальное приближение подбирается индивидуально.
            \end{itemize}
            \textbf{Класс проектных задач}
            \begin{itemize}
                \item меняются даты старта и параметры конечной орбиты;
                \item возникают разные ветви решений;
                \item требуется согласованный выбор начальных сопряжённых переменных.
            \end{itemize}
        \end{column}
    \end{columns}
\end{frame}

\begin{frame}
    \frametitle{Степень разработанности и место работы}
    \begin{columns}[T,onlytextwidth]
        \begin{column}{0.31\textwidth}
            \textbf{Аналитические методы}
            \begin{itemize}
                \item конечные формулы;
                \item специальные постановки;
                \item ограниченные диапазоны применимости.
            \end{itemize}
        \end{column}
        \begin{column}{0.31\textwidth}
            \textbf{Численные методы}
            \begin{itemize}
                \item решение конкретных краевых задач;
                \item высокая чувствительность к начальному приближению;
                \item ветвление решений.
            \end{itemize}
        \end{column}
        \begin{column}{0.31\textwidth}
            \textbf{Задача диссертации}
            \begin{itemize}
                \item построить модельное многообразие;
                \item получить конечные формулы;
                \item связать их с начальным приближением для прикладной КЗ.
            \end{itemize}
        \end{column}
    \end{columns}
\end{frame}

\begin{frame}
    \frametitle{Цель и задачи}
    \textbf{Цель:} разработать и апробировать численно-аналитический метод построения начальных приближений для энергетически оптимальных межпланетных перелётов с непрерывной тягой на основе аппроксимации модельного многообразия решений задачи встречи.

    \medskip
    \begin{enumerate}
        \item разработать методику построения и аппроксимации модельного многообразия решений;
        \item исследовать плоскую круговую модель и получить конечные формулы;
        \item разработать линейные преобразования для эллиптичной и пространственной конечной орбиты;
        \item построить и апробировать алгоритм выбора начального приближения.
    \end{enumerate}
\end{frame}

\begin{frame}
    \frametitle{Научная новизна}
    \begin{columns}[T,onlytextwidth]
        \begin{column}{0.48\textwidth}
            \begin{block}{Модельное многообразие}
                Предложена методика построения и аппроксимации регулярных ветвей решений задачи встречи.
            \end{block}
            \begin{block}{Конечные формулы}
                Получены аппроксимации продолжительности, фазы, функционала и начальных сопряжённых переменных.
            \end{block}
        \end{column}
        \begin{column}{0.48\textwidth}
            \begin{block}{\(R\)-преобразования}
                Разработаны линейные преобразования начальных значений сопряжённых переменных.
            \end{block}
            \begin{block}{Алгоритм}
                Построен алгоритм выбора начального приближения и выполнена его программная реализация.
            \end{block}
        \end{column}
    \end{columns}
\end{frame}

\section{Глава 1. Методика}

\begin{frame}
    \frametitle{Глава 1. Методика построения и аппроксимации}
    \begin{columns}[T,onlytextwidth]
        \begin{column}{0.46\textwidth}
            \begin{enumerate}
                \item Энергетически оптимальная модельная задача встречи.
                \item Построение многообразия дискретным продолжением.
                \item Аппроксимация характеристик с помощью символьной регрессии.
            \end{enumerate}
        \end{column}
        \begin{column}{0.5\textwidth}
            \centering
            \includegraphics[width=0.95\linewidth]{methodology_cycle}
        \end{column}
    \end{columns}
\end{frame}

\begin{frame}
    \frametitle{Идеально регулируемая тяга как модельная КЗ}
    \begin{columns}[T,onlytextwidth]
        \begin{column}{0.48\textwidth}
            \small
            \[
                J=\int_{t_0}^{t_f}\|\mathbf a(t)\|^2\,dt
            \]
            \[
                \mathbf y_0=
                \left(
                \mathbf r_0,\mathbf v_0,
                \mathbf p_{r0},\mathbf p_{v0}
                \right)
            \]
            \begin{itemize}
                \item постановка сохраняет двухточечную краевую задачу;
                \item управление определяется начальными значениями сопряжённых переменных;
                \item качество приближения оценивается через последующее уточнение КЗ.
            \end{itemize}
        \end{column}
        \begin{column}{0.48\textwidth}
            \centering
            \includegraphics[height=0.68\textheight]{virtual_planet}
        \end{column}
    \end{columns}
\end{frame}

\begin{frame}
    \frametitle{Параметры модельной задачи}
    \begin{columns}[T,onlytextwidth]
        \begin{column}{0.5\textwidth}
            \centering
            \includegraphics[width=\linewidth]{virtual_planet}
        \end{column}
        \begin{column}{0.46\textwidth}
            \begin{itemize}
                \item начальная конфигурация задаётся синодической фазой;
                \item конечная точка встречи задаётся угловой дальностью;
                \item отношение радиусов круговых орбит задаёт параметр \(\kappa\);
                \item витковая дальность используется как параметр \(\lambda\).
            \end{itemize}
            \[
                (\lambda,\kappa)
                \longrightarrow
                T,\ \Psi,\ J,\ \mathbf z_0
            \]
        \end{column}
    \end{columns}
\end{frame}

\begin{frame}
    \frametitle{Построение многообразия дискретным продолжением}
    \begin{columns}[T,onlytextwidth]
        \begin{column}{0.58\textwidth}
            \centering
            \includegraphics[width=\linewidth]{continuation_paths}
        \end{column}
        \begin{column}{0.38\textwidth}
            \begin{itemize}
                \item базовое решение продолжается малыми шагами по пространству параметров;
                \item соседние решения используются как начальные приближения;
                \item вычислительные барьеры выделяют границы регулярной области.
            \end{itemize}
        \end{column}
    \end{columns}
\end{frame}

\begin{frame}
    \frametitle{Аппроксимация символьной регрессией}
    \begin{columns}[T,onlytextwidth]
        \begin{column}{0.48\textwidth}
            \[
            \begin{array}{c}
            \text{многообразие решений}\\
            \Downarrow\\
            \text{выбор структуры формул}\\
            \Downarrow\\
            \text{оценка коэффициентов}\\
            \Downarrow\\
            \text{конечные формулы}
            \end{array}
            \]
        \end{column}
        \begin{column}{0.48\textwidth}
            \begin{itemize}
                \item структура формул выбирается автоматически;
                \item проверяются отсутствие сингулярностей и точность на отложенной выборке;
                \item учитываются предельные случаи, допускающие аналитическую интерпретацию.
            \end{itemize}
        \end{column}
    \end{columns}
\end{frame}

\begin{frame}
    \frametitle{Результаты главы 1}
    \begin{enumerate}
        \item Сформулирована регулярная модельная краевая задача встречи с идеально регулируемой тягой.
        \item Разработан способ построения многообразия решений методом дискретного продолжения.
        \item Задана методика получения конечных формул для характеристик экстремалей и начальных сопряжённых переменных.
    \end{enumerate}
\end{frame}

\section{Глава 2. Плоская круговая модель}

\begin{frame}
    \frametitle{Глава 2. Плоская круговая модель}
    \begin{columns}[T,onlytextwidth]
        \begin{column}{0.42\textwidth}
            \begin{enumerate}
                \item Семейства решений и ветвление.
                \item Кривая локальных оптимумов.
                \item Конечные формулы для \(T\), \(\Psi\), \(J\), \(\mathbf z_0\).
                \item Верификация качества аппроксимаций.
            \end{enumerate}
        \end{column}
        \begin{column}{0.54\textwidth}
            \centering
            \includegraphics[width=\linewidth]{5families}
        \end{column}
    \end{columns}
\end{frame}

\begin{frame}
    \frametitle{Семейства решений и выбор ветви}
    \begin{columns}[T,onlytextwidth]
        \begin{column}{0.62\textwidth}
            \centering
            \includegraphics[width=\linewidth]{5families}
        \end{column}
        \begin{column}{0.34\textwidth}
            \begin{itemize}
                \item многообразие состоит из нескольких регулярных семейств;
                \item разные семейства могут соответствовать близким датам и дальностям;
                \item выбор ветви определяет начальное приближение для непрямого метода.
            \end{itemize}
        \end{column}
    \end{columns}
\end{frame}

\begin{frame}
    \frametitle{Кривая локальных оптимумов и точки оптимальной дальности}
    \begin{columns}[T,onlytextwidth]
        \begin{column}{0.6\textwidth}
            \centering
            \includegraphics[width=\linewidth]{ParetoToF}
        \end{column}
        \begin{column}{0.36\textwidth}
            \begin{itemize}
                \item КЛО выделяет регулярную структуру внутри многообразия;
                \item для фиксированной начальной фазы выбирается точка оптимальной дальности;
                \item эта точка задаёт ветвь, время перелёта и начальные сопряжённые переменные.
            \end{itemize}
        \end{column}
    \end{columns}
\end{frame}

\begin{frame}
    \frametitle{Конечная формула для продолжительности перелёта}
    \begin{columns}[T,onlytextwidth]
        \begin{column}{0.48\textwidth}
            \scriptsize
            \[
            \widehat{T}(\lambda^\star,\kappa)=
            \frac{2}{27}\kappa\ln\kappa+
            \lambda^\star
            \left(\kappa-\frac{\sqrt6}{9}\kappa\ln\kappa\right)
            \]
            \[
            \widehat{T}(\lambda^\star,1)=\lambda^\star
            \]
            \begin{itemize}
                \item формула согласована с предельным случаем совпадающих круговых орбит;
                \item время перелёта становится явной функцией параметров многообразия.
            \end{itemize}
        \end{column}
        \begin{column}{0.48\textwidth}
            \centering
            \includegraphics[width=\linewidth]{main_T_isolines}
        \end{column}
    \end{columns}
\end{frame}

\begin{frame}
    \frametitle{Формулы для фазы и функционала}
    \begin{columns}[T,onlytextwidth]
        \begin{column}{0.47\textwidth}
            \footnotesize
            \[
            \widehat{\Psi}(\lambda^\star,\kappa)=
            \mathcal A_0^{(\Psi)}(\kappa)+
            \mathcal A_1^{(\Psi)}(\kappa)\lambda^\star
            \]
            \[
            \begin{aligned}
            \widehat J &=
            \frac{(C_0+C_1/\kappa+C_2/\kappa^2)\ln\kappa}{D},\\
            D&=\lambda^\star+
            C_3(\sin(2\pi\lambda^\star+C_4-1/\lambda^\star)
            -\lambda^\star)e^{-\lambda^\star}
            \end{aligned}
            \]
        \end{column}
        \begin{column}{0.49\textwidth}
            \centering
            \includegraphics[width=0.49\linewidth]{main_phase_isolines}
            \includegraphics[width=0.49\linewidth]{main_J_isolines}

            \vspace{0.1cm}
            \footnotesize Аппроксимации задают скалярные слои модельного многообразия.
        \end{column}
    \end{columns}
\end{frame}

\begin{frame}
    \frametitle{Формулы для начальных сопряжённых переменных}
    \begin{columns}[T,onlytextwidth]
        \begin{column}{0.36\textwidth}
            \[
            \widehat{\mathbf z}^{\mathrm{pl}}_0=
            \left(
            \widehat p_{rx},
            \widehat p_{ry},
            \widehat p_{vx},
            \widehat p_{vy}
            \right)
            \]
            \begin{itemize}
                \item полученный вектор используется как вход в расширенную систему;
                \item коэффициенты формул приведены в приложениях диссертации.
            \end{itemize}
        \end{column}
        \begin{column}{0.6\textwidth}
            \centering
            \includegraphics[width=\linewidth]{main_p_isolines}
        \end{column}
    \end{columns}
\end{frame}

\begin{frame}
    \frametitle{Качество конечных формул}
    \begin{columns}[T,onlytextwidth]
        \begin{column}{0.6\textwidth}
            \centering
            \includegraphics[height=0.68\textheight]{main_T_error_heatmap}

            \vspace{0.1cm}
            \footnotesize Ошибка аппроксимации времени перелёта
        \end{column}
        \begin{column}{0.36\textwidth}
            \begin{itemize}
                \item качество оценивается относительно задачи начального приближения;
                \item формулы должны обеспечить вход в область сходимости непрямого метода;
                \item окончательное решение получается последующим численным уточнением.
            \end{itemize}
        \end{column}
    \end{columns}
\end{frame}

\begin{frame}
    \frametitle{Результаты главы 2}
    \begin{enumerate}
        \item Показано существование точек оптимальной дальности на регулярных ветвях многообразия.
        \item Получены конечные формулы для продолжительности перелёта, синодической фазы, функционала и начальных сопряжённых переменных.
        \item Построено аппроксимированное многообразие решений плоской круговой задачи встречи.
    \end{enumerate}
\end{frame}

\section{Глава 3. R-преобразования}

\begin{frame}
    \frametitle{Глава 3. Линейные \(R\)-преобразования}
    \begin{columns}[T,onlytextwidth]
        \begin{column}{0.42\textwidth}
            \begin{enumerate}
                \item Общий смысл перехода между краевыми задачами.
                \item Преобразование начального фазового вектора.
                \item Переход к эллиптичной конечной орбите.
                \item Переход к пространственной конечной орбите.
            \end{enumerate}
        \end{column}
        \begin{column}{0.54\textwidth}
            \centering
            \includegraphics[width=\linewidth]{R_small_change_demo}
        \end{column}
    \end{columns}
\end{frame}

\begin{frame}
    \frametitle{Смысл линейного \(R\)-преобразования}
    \begin{columns}[T,onlytextwidth]
        \begin{column}{0.42\textwidth}
            \[
            \begin{array}{c}
            \mathrm{КЗ}_1\rightarrow\mathrm{экстремаль}_1\\
            \xrightarrow{R}\\
            \mathrm{экстремаль}_2\rightarrow\mathrm{КЗ}_2
            \end{array}
            \]
            \begin{itemize}
                \item преобразуется начальный вектор расширенной системы;
                \item меняются параметры краевой задачи;
                \item полученная экстремаль используется как начальное приближение для новой постановки.
            \end{itemize}
        \end{column}
        \begin{column}{0.54\textwidth}
            \centering
            \includegraphics[width=\linewidth]{R_small_change_demo}
        \end{column}
    \end{columns}
\end{frame}

\begin{frame}
    \frametitle{\(R\)-преобразование I типа}
    \begin{columns}[T,onlytextwidth]
        \begin{column}{0.43\textwidth}
            \begin{itemize}
                \item задаёт переход к новому начальному фазовому вектору;
                \item согласованно преобразует начальные сопряжённые переменные;
                \item отличается от простого поворота вектора сопряжённых переменных.
            \end{itemize}
        \end{column}
        \begin{column}{0.53\textwidth}
            \centering
            \includegraphics[width=\linewidth]{R_type_I_demo}
        \end{column}
    \end{columns}
\end{frame}

\begin{frame}
    \frametitle{\(R\)-преобразование II типа}
    \begin{columns}[T,onlytextwidth]
        \begin{column}{0.38\textwidth}
            \[
            (e,\omega)
            \longrightarrow
            (\zeta,\chi)
            \longrightarrow
            R_{e\omega}
            \]
            \begin{itemize}
                \item задаются эксцентриситет и аргумент перицентра;
                \item учитывается деформация витковой дальности;
                \item переход используется при построении прикладной конечной орбиты.
            \end{itemize}
        \end{column}
        \begin{column}{0.58\textwidth}
            \centering
            \includegraphics[width=\linewidth]{R_type_II_eomega_example}
        \end{column}
    \end{columns}
\end{frame}

\begin{frame}
    \frametitle{\(R\)-преобразование III типа}
    \begin{columns}[T,onlytextwidth]
        \begin{column}{0.38\textwidth}
            \[
            (i,\Omega)
            \longrightarrow
            (j,\gamma)
            \longrightarrow
            R_{i\Omega}
            \]
            \begin{itemize}
                \item задаются наклонение и долгота восходящего узла;
                \item плоская модель переводится в пространственную постановку;
                \item преобразование сохраняет вычислительную роль начального приближения.
            \end{itemize}
        \end{column}
        \begin{column}{0.58\textwidth}
            \centering
            \includegraphics[width=\linewidth]{R_type_III_iOmega_example}
        \end{column}
    \end{columns}
\end{frame}

\begin{frame}
    \frametitle{Композиция \(R\)-преобразований}
    \begin{columns}[T,onlytextwidth]
        \begin{column}{0.38\textwidth}
            \[
            R_{\Sigma}=R_{i\Omega}R_{e\omega}R_0
            \]
            \begin{itemize}
                \item сначала задаётся начальный фазовый вектор;
                \item затем форма конечной орбиты;
                \item затем пространственная ориентация.
            \end{itemize}
        \end{column}
        \begin{column}{0.58\textwidth}
            \centering
            \includegraphics[height=0.7\textheight]{R_composition_sequential_example}
        \end{column}
    \end{columns}
\end{frame}

\begin{frame}
    \frametitle{Результаты главы 3}
    \begin{enumerate}
        \item Предложено линейное преобразование специального вида для начальных значений сопряжённых переменных.
        \item Получены конечные формулы, связывающие управляющие параметры преобразований с элементами конечной орбиты.
        \item Построена композиция преобразований для перехода к эллиптичной и пространственной конечной орбите.
    \end{enumerate}
\end{frame}

\section{Глава 4. Применение}

\begin{frame}
    \frametitle{Глава 4. Применение в баллистическом проектировании}
    \begin{columns}[T,onlytextwidth]
        \begin{column}{0.42\textwidth}
            \begin{enumerate}
                \item Управляющие параметры начального приближения.
                \item Алгоритм построения \(\mathbf y_0^{init}\).
                \item Проверка на перелётах Земля--Марс и Земля--Венера.
                \item Программная реализация в ILTT Toolbox.
            \end{enumerate}
        \end{column}
        \begin{column}{0.54\textwidth}
            \centering
            \includegraphics[width=\linewidth]{Venus_initial_guess_after_transfer}
        \end{column}
    \end{columns}
\end{frame}

\begin{frame}
    \frametitle{Управляющие параметры и алгоритм}
    \small
    \[
    t_0,\ \mathrm{target\ orbit}
    \rightarrow
    \Phi,\ (a,e,\omega,i,\Omega)
    \rightarrow
    \boldsymbol q
    \rightarrow
    \widehat{\mathbf z}_0^{\mathrm{pl}}
    \xrightarrow{R_{\Sigma}}
    \mathbf y_0^{\mathrm{init}}
    \rightarrow
    \mathrm{BVP}
    \]
    \begin{center}
        \footnotesize
        \begin{tabular}{@{}lll@{}}
            \toprule
            Параметры & Роль & Источник \\
            \midrule
            \(\lambda,\kappa\) & точка на модельном многообразии & глава 2 \\
            \(\zeta,\chi\) & форма конечной орбиты & \(R\)-II \\
            \(j,\gamma\) & пространственная ориентация & \(R\)-III \\
            \bottomrule
        \end{tabular}
    \end{center}
    Вектор \(\boldsymbol q\) является вычислительным интерфейсом между конечными формулами, \(R\)-преобразованиями и численным уточнением краевой задачи.
\end{frame}

\begin{frame}
    \frametitle{Пример: перелёт Земля--Марс}
    \begin{columns}[T,onlytextwidth]
        \begin{column}{0.64\textwidth}
            \centering
            \includegraphics[width=\linewidth]{chap4_mars_cleanver_rebuilt_Oxy}
        \end{column}
        \begin{column}{0.32\textwidth}
            \begin{itemize}
                \item дата старта и целевая орбита задают входные параметры;
                \item конечные формулы дают оценку \(T\) и начальные сопряжённые переменные;
                \item после \(R\)-преобразований выполняется уточнение краевой задачи.
            \end{itemize}
        \end{column}
    \end{columns}
\end{frame}

\begin{frame}
    \frametitle{Качество начального приближения}
    \begin{columns}[T,onlytextwidth]
        \begin{column}{0.48\textwidth}
            \begin{itemize}
                \item Проверка выполнена для множеств перелётов Земля--Марс и Земля--Венера.
                \item Для Марса рассматривался период великих противостояний порядка 17 лет.
                \item Для Венеры рассматривался характерный восьмилетний цикл.
                \item Цель проверки --- оценить невязку начального приближения и сходимость уточнения.
            \end{itemize}
        \end{column}
        \begin{column}{0.48\textwidth}
            \centering
            \footnotesize
            \begin{tabular}{@{}lcc@{}}
                \toprule
                Параметр & Земля--Марс & Земля--Венера \\
                \midrule
                \(e\) & \(0.003\) & \(0.02\) \\
                \(i\) & \(0.004^\circ\) & \(0.2^\circ\) \\
                \(\Omega\) & \(4^\circ\) & \(2^\circ\) \\
                \(\omega\) & \(11^\circ\) & \(26^\circ\) \\
                \(\varphi\) & \(3^\circ\) & \(0.7^\circ\) \\
                \(a\) & \(0.06\) а.е. & \(0.004\) а.е. \\
                \(T\) & \(8.5\) сут. & \(0.7\) сут. \\
                \bottomrule
            \end{tabular}
        \end{column}
    \end{columns}
\end{frame}

\begin{frame}
    \frametitle{ILTT Toolbox}
    \begin{columns}[T,onlytextwidth]
        \begin{column}{0.48\textwidth}
            \begin{itemize}
                \item Результаты реализованы в программном инструменте \texttt{ILTT Toolbox}.
                \item Пользователь задаёт дату старта и параметры целевой орбиты.
                \item На выходе формируется начальное приближение для последующего решения краевой задачи.
                \item Получено свидетельство о государственной регистрации программы для ЭВМ.
            \end{itemize}
        \end{column}
        \begin{column}{0.48\textwidth}
            \centering
            \footnotesize
            \begin{tabular}{@{}p{0.25\linewidth}p{0.65\linewidth}@{}}
                \toprule
                Вход & дата старта, целевая орбита \\
                Метод & конечные формулы и \(R_\Sigma\) \\
                Выход & \(\mathbf y_0^{init}\), \(\widehat T\) \\
                Проверка & численное уточнение КЗ \\
                \bottomrule
            \end{tabular}
        \end{column}
    \end{columns}
\end{frame}

\begin{frame}
    \frametitle{Переход к прикладной расчётной цепочке}
    \begin{columns}[T,onlytextwidth]
        \begin{column}{0.58\textwidth}
            \centering
            \includegraphics[width=\linewidth]{Venus_deltaV_by_start_date}
        \end{column}
        \begin{column}{0.38\textwidth}
            \small
            \textbf{Продолжение расчёта}
            \begin{itemize}
                \item ИРТ-приближение используется как вход;
                \item затем учитываются \(v_\infty\) и постоянная скорость истечения;
                \item результат применяется при анализе окон старта.
            \end{itemize}
        \end{column}
    \end{columns}
\end{frame}
